% !TEX program = xelatex
\documentclass{article}
\usepackage{geometry}\geometry{margin=23mm}
\usepackage{fontspec}\setmainfont{Times New Roman}
\usepackage[ngerman]{babel}
\begin{document}
\title{Das Messbrett}\author{}\date{14. Mai 2026}\maketitle
\begin{abstract}Zwei Marken zeigen die ursprüngliche Lage des Brettes.\end{abstract}
\tableofcontents
\section{Die Marken}\begin{quote}„Zuerst die untere Marke ablesen.“\end{quote}
\begin{figure}[ht]\centering\fbox{\rule{0pt}{8mm}\rule{33mm}{0pt}}\caption{Lage des Brettes}\end{figure}
\section{Höhe}
\begin{table}[ht]\centering\begin{tabular}{cr}Marke&Höhe (mm)\\\hline 1&12\\2&18\end{tabular}\caption{Höhe der Marken}\end{table}
\begin{thebibliography}{1}\bibitem{record}\emph{Werkstattheft}, 2025, \chaptername~2.\end{thebibliography}
\end{document}
